% context: 9th-grade physics homework, Unit 2 kinematics, after the free-fall lesson
% topic: mental math, no calculator, answers only — metric length prefixes, squares
%        and square roots, (1/2)at^2 with round numbers, m/s and km/h, US units,
%        fractions to decimals, fractions of an hour, a falling stone with g = 10 m/s^2
% structure: two pages; page 1 prefixes, squares, and speed units; page 2 fractions,
%        time, and the falling stone. Follows 2-1DN_Mental-math with new numbers
% data: Empire State Building roof 381 m (contexts.md, Distance); Giancoli problem 40
%        drops King Kong from it
\documentclass[12pt, twoside]{article}
%\documentclass[12pt, twoside]{article}
\usepackage[letterpaper, margin=1in, headsep=0.2in]{geometry}
\setlength{\headheight}{0.6in}
%\usepackage[english]{babel}
\usepackage[utf8]{inputenc}
\usepackage{microtype}
\usepackage{amsmath}
\usepackage{amssymb}
%\usepackage{amsfonts}
\usepackage[nomessages]{fp} %\FPeval{\var-name}{2*sin(pi/6)}
\usepackage{siunitx} %units in math. eg 20\milli\meter
\usepackage{yhmath} % for arcs, overparenth command
\usepackage{tikz} %graphics
\usetikzlibrary{quotes, angles, arrows, arrows.meta}
\usepackage{graphicx} %consider setting \graphicspath{{images/}}
\usepackage{parskip} %no paragraph indent
\usepackage{enumitem}
\usepackage{multicol}
\usepackage{venndiagram}

\usepackage{fancyhdr}
\pagestyle{fancy}
\fancyhf{}
\renewcommand{\headrulewidth}{0pt} % disable the underline of the header
\raggedbottom
\hfuzz=2mm %suppresses overfull box warnings

\usepackage{hyperref}
\usepackage{float}

\newcommand{\blank}[1][2.4cm]{\rule{#1}{0.4pt}}

\title{Physics}
\author{Chris Huson}
\date{October 2026}

\fancyhead[LE]{\thepage}
\fancyhead[RO]{\thepage \\ Nickname: \hspace{2.25cm} \,\\ \,}
\fancyhead[LO]{La Scuola d'Italia / Huson / Physics \\* 30 October 2026}

\begin{document}

\subsubsection*{2.5 Homework: Mental Math (no calculator)}

Work each question in your head --- no calculator, and no need to show work. Write
just the answer, with its units.

\vspace{0.15cm}

\textbf{Reference.} \quad
kilo (k) = 1000 \quad centi (c) = $\dfrac{1}{100}$ \quad
milli (m) = $\dfrac{1}{1000}$ \quad 1 m/s = 3.6 km/h

\vspace{0.2cm}

\begin{enumerate}[itemsep=0.4cm]

\item Fill in the blanks.
  \begin{multicols}{2}
    \begin{enumerate}[itemsep=0.5cm]
      \item 1.6 km = \blank\ m
      \item 85 cm = \blank\ m
      \item 0.25 m = \blank\ cm
      \item 4 mm = \blank\ cm
      \item 2.2 m = \blank\ mm
      \item 500 mm = \blank\ m
    \end{enumerate}
  \end{multicols}

\item Squares and square roots.
  \begin{multicols}{4}
    \begin{enumerate}[itemsep=0.6cm]
      \item $3^2 = $ \blank[1.1cm]
      \item $7^2 = $ \blank[1.1cm]
      \item $12^2 = $ \blank[1.1cm]
      \item $20^2 = $ \blank[1.1cm]
      \item $0.5^2 = $ \blank[1.1cm]
      \item $0.3^2 = $ \blank[1.1cm]
      \item $\sqrt{81} = $ \blank[1.1cm]
      \item $\sqrt{144} = $ \blank[1.1cm]
    \end{enumerate}
  \end{multicols}

\item Work out $\frac{1}{2} \times a \times t^2$. Square first, then multiply.
  \begin{multicols}{2}
    \begin{enumerate}[itemsep=0.5cm]
      \item $\frac{1}{2} \times 10 \times 2^2 = $ \blank[1.8cm]
      \item $\frac{1}{2} \times 10 \times 3^2 = $ \blank[1.8cm]
      \item $\frac{1}{2} \times 4 \times 5^2 = $ \blank[1.8cm]
      \item $\frac{1}{2} \times 0.2 \times 10^2 = $ \blank[1.8cm]
    \end{enumerate}
  \end{multicols}

\item Use 1 m/s = 3.6 km/h.
  \begin{multicols}{2}
    \begin{enumerate}[itemsep=0.5cm]
      \item 10 m/s = \blank[1.8cm] km/h
      \item 20 m/s = \blank[1.8cm] km/h
      \item 108 km/h = \blank[1.8cm] m/s
      \item 360 km/h = \blank[1.8cm] m/s
    \end{enumerate}
  \end{multicols}

\item Circle the best answer.
  \begin{enumerate}[itemsep=0.5cm]
    \item One foot is about how many meters? \hfill
      (A) 0.3 \quad\; (B) 0.5 \quad\; (C) 3.28 \quad\; (D) 12
    \item One mile is about how many km? \hfill
      (A) 0.6 \quad\; (B) 1.6 \quad\; (C) 2.54 \quad\; (D) 5280
    \item 60 mi/h is about how many m/s? \hfill
      (A) 16 \quad\; (B) 27 \quad\; (C) 60 \quad\; (D) 97
  \end{enumerate}

\newpage

\item Write each fraction as a decimal.
  \begin{multicols}{4}
    \begin{enumerate}[itemsep=0.6cm]
      \item $\dfrac{2}{5} = $ \blank[1.6cm]
      \item $\dfrac{4}{5} = $ \blank[1.6cm]
      \item $\dfrac{3}{8} = $ \blank[1.6cm]
      \item $\dfrac{7}{8} = $ \blank[1.6cm]
      \item $\dfrac{1}{20} = $ \blank[1.6cm]
      \item $\dfrac{1}{25} = $ \blank[1.6cm]
      \item $\dfrac{1}{6} = $ \blank[1.6cm]
      \item $\dfrac{5}{6} = $ \blank[1.6cm]
    \end{enumerate}
  \end{multicols}

\item Fill in the blanks.
  \begin{multicols}{2}
    \begin{enumerate}[itemsep=0.6cm]
      \item 0.25 hour = \blank[1.8cm] min
      \item 0.3 hour = \blank[1.8cm] min
      \item 1.5 min = \blank[1.8cm] s
      \item 150 s = \blank[1.8cm] min
      \item 12 min = \blank[1.8cm] h (decimal)
      \item 20 min = \blank[1.8cm] h (fraction)
      \item 1 day = \blank[1.8cm] min
      \item 1.75 hours = \blank[1.2cm] h \blank[1.2cm] min
    \end{enumerate}
  \end{multicols}

\item Physics meets twice a week for 45 minutes.
  \begin{enumerate}[itemsep=0.6cm]
    \item How many hours of physics is that each week? \hfill \blank
    \item A school year has about 36 weeks of classes. About how many hours of
      physics is that? \hfill \blank
  \end{enumerate}

\item A 10K race is 10 kilometers long. About how many miles is that? \hfill \blank

\item A stone is dropped from a tall building. Use $g \approx 10$ m/s$^2$: each second
  the stone gains 10 m/s of speed, and after $t$ seconds it has fallen
  $\frac{1}{2} \times 10 \times t^2$ meters.
  \begin{enumerate}[itemsep=0.6cm]
    \item Fill in the table.

      \vspace{0.1cm}
      \renewcommand{\arraystretch}{1.5}
      \begin{tabular}{| c | p{1.2cm} | p{1.2cm} | p{1.2cm} | p{1.2cm} | p{1.2cm} |}
        \hline
        $t$ (s) & \centering 1 & \centering 2 & \centering 3 &
          \centering 4 & \centering\arraybackslash 5 \\ \hline
        speed (m/s) & & & & & \\ \hline
        distance fallen (m) & & & & & \\ \hline
      \end{tabular}
    \item How many seconds does it take to fall 320 m? \hfill \blank
    \item The roof of the Empire State Building is 381 m high. A stone dropped from
      it lands between which two whole numbers of seconds? \hfill
      \blank[1.2cm] and \blank[1.2cm] s
  \end{enumerate}

\end{enumerate}
\end{document}
